%% ===========================================================================
\section{Discussion}
\label{sec:discussion}
%% ===========================================================================

This section and the next are the drafting tool's interpretation of the results, written so that the study can be read as a whole. They are one reading of the numbers, not the author's, and the paper's discussion is to be written independently.

\subsection{What sets the signal a filter removes}

The results are consistent with the hypothesis of the proposal: the fraction of a transient that a filter removes depends mainly on how much of the filter's input lies inside the transient. For a linear filter, Eq.~\eqref{eq:retention} gives $\rho = 1 - \langle P(s), s\rangle / \langle s, s \rangle$, so the loss is the part of the transient that the filter's own prediction reproduces. For the far stack at $R = 0$, the prediction is the network mean of $s$, and a Gaussian transient contributes to that mean in proportion to the share of stations it covers. Computing this share from the station positions alone, ignoring data gaps, gives median retentions of 1.00, 0.985, 0.95 and 0.85 for widths of 10, 25, 50 and 100~km, against measured values of 1.00, 0.98, 0.94 and 0.83. The far stack therefore loses a 100-km transient for a simple geometric reason: such a transient covers a noticeable fraction of a network that is itself only a few hundred kilometres across.

For the nearest-neighbour filters the same reasoning works the other way. The 16th neighbour lies at a median of 40~km, so a 25-km transient covers most of a station's neighbours and a 100-km transient covers all of them; M1 keeps 0.25 and 0.00 of these. Increasing the exclusion radius moves the neighbours outside the transient, and retention rises when $R$ becomes comparable with $2L$ to $3L$: M1 keeps 0.83 of 25-km transients at $R = 50$~km, 0.73 of 50-km transients at $R = 100$~km and 0.70 of 100-km transients at $R = 200$~km. The price is that the remaining neighbours are farther away and fewer stations share their noise with the target, so noise removed falls from 14.1\% to 4.7\%, and spurious signal rises, because a transient now reaches stations whose neighbours are inside it while they themselves are not.

The exclusion radius is therefore not a parameter with an optimum but a choice of which transients to protect. For the far stack, raising $R$ from 0 to 200~km raises retention of 100-km transients from 0.83 to 0.96 at a cost of 0.6 points of noise removed and 0.006 in reconstruction error; at 400~km, all transients are kept, but noise removed falls by 2.2 points and spurious signal doubles. Which point is preferable depends on the transients a study is looking for.

\subsection{How many stations, and how far}

On the validation years the weighted mean of the 256 nearest stations and the far stack had the same reconstruction error, which suggested that the number of stations averaged, not their distance, is what reduces the error. The test years do not support that reading in its simple form: the far stack's error is lower by 0.04, and the 64-neighbour mean has a lower error than the 256-neighbour one. The two filters differ in more than the number of stations, since M1 weights its neighbours by distance with a length scale tuned on the validation years, and the comparison cannot separate the two effects. What holds in both periods is that averaging more stations keeps more signal (25-km retention 0.25, 0.53, 0.84 and 0.98 for 16, 64, 256 and all stations), and that no filter over a subset of nearby stations removed more noise than the mean over all of them on the test years.

\subsection{Local filters and the far field}

The fault-slip events show a pattern that the Gaussian transients cannot: the nearest-neighbour filters keep afterslip close to the fault and remove it farther away. Within 5~km of the patch, every filter keeps 0.94--1.01; at 15--40~km, M1 keeps 0.11 and M3k 0.13, while the far stack keeps 1.00. Near the fault the displacement varies over distances shorter than the spacing of the neighbours, and on a strike-slip fault the two sides move in opposite directions, so the neighbours' mean does not reproduce the motion at the target. Far from the fault the displacement field is smooth over the scale of the neighbourhood, and a neighbour average reproduces it, as it does for a wide Gaussian transient. The far field of afterslip and slow slip constrains the depth and moment of the slip, so a local filter removes the part of the signal that a source inversion relies on. The Gaussian transients are in this sense pessimistic near faults and optimistic far from them, and the whole-event retention, dominated by the near-fault stations with the largest motion (0.79--1.00 for every filter), hides the difference.

\subsection{Learned filters}

Learning helps a filter that is restricted to the same inputs: the learned kernel has a lower error than the fixed Gaussian kernel over 16 neighbours (0.595 against 0.621) and over 256 (0.562 against 0.606), and removes more noise, at the same or lower retention. None of the learned filters, however, has a reconstruction error distinguishable from that of the far stack on the test years, on seen or on held-out stations, and every one keeps less planted signal. The attention network is built as the far stack plus a learned correction; without planted transients in training, the correction removes part of every coherent transient (100-km retention 0.59 against 0.83) without a measurable reduction in error, and with them (M5-aug), it keeps signal almost as well as the far stack and removes the same noise. On the validation years, the learned kernel over 256 neighbours had a lower error than the far stack with an interval excluding zero; that this advantage disappeared on the test years suggests that part of what was learned was specific to the validation period, or that the test years, which contain a different mix of noise and earthquakes, do not reward it. On this network and with daily data, the far stack is the baseline the learned filters did not beat, and the measure that separates them from it is signal kept rather than noise removed.

\subsection{Target-inclusive filters and the measures}

The target-inclusive filters illustrate why scatter reduction alone misleads. M2-self is M2 with the target added to the median: it removes 19.3\% of the scatter against 12.2\% for M2, while keeping the same signal (0.24 against 0.28 at 25~km) and real offsets (0.14 against 0.12). The extra noise removed is part of the station's own noise, which a filter that reads the target can subtract and a leave-one-out filter cannot; it says nothing about the common mode. C1, the network mean with the target included, behaves like the far stack at $R = 0$ on every measure, since one station out of 1,251 contributes little to the mean. C3 shows a different cost: with one region and one component its spurious signal is 1.58 times the planted amplitude for 100-km transients, because a single component fitted to the whole network spreads part of any large transient over every station; more regions reduce the spurious signal but remove more of the transient. Velocity repeatability rewards the same behaviour as scatter reduction: M2-self has the smallest vertical $|\Delta v|$, and a filter that absorbs real changes in velocity improves this measure. Reconstruction error did not predict retention either: M3k has a lower error than M1 and M2 and keeps as little. Retention and spurious signal are therefore not refinements of the noise measures but independent axes, and a comparison of filters that omits them can rank a filter that removes signal above one that does not.

\subsection{Relation to Bachelot et al.}

The far stack at 400~km (D1) keeps every planted transient and all of the afterslip, consistent with the choice of \citet{bachelot2025} to exclude nearby stations, and the far stack performing as well as the learned filters here parallels their finding that a stack of distant stations performed comparably to their network. The audit adds what the fixed distance costs: compared with the far stack at $R = 0$, D1 removes 2.2 points less of the scatter and creates twice the spurious signal for 100-km transients, while the gain in retention is confined to transients wider than about 50~km. On this network, a smaller radius keeps transients up to 50~km almost completely (0.96 at $R = 50$~km).

\subsection{Limitations}

The audit uses one dense network in a strike-slip setting; in a sparse network, such as Türkiye's, where the median distance to the nearest NGL station is about 60~km, every filter's neighbours are farther away and the trade-off shifts. The fault-slip events use an elastic half-space, uniform tapered slip and two templates. The real offsets have no ground truth, and on the test years two thirds of them come from one earthquake. Only daily solutions are used. The test years were looked at while the scoring rules were revised, and several comparisons differ between the validation and test years, so the prospective evaluation is needed before any of the smaller differences reported here can be relied on. Velocity repeatability in the north component worsens with almost every filter for reasons that have not been examined.

%% ===========================================================================
\section{Conclusions}
\label{sec:conclusions}
%% ===========================================================================

This section is the drafting tool's summary of its own interpretation (Section~\ref{sec:discussion}).

\begin{enumerate}
  \item The signal a common-mode filter removes is set by how much of its input lies inside the transient. Filters built on the 16 nearest stations remove most of any transient wider than about 25~km; the mean of all stations loses only the share of a transient that enters the network mean, 0.17 of a 100-km transient on this network.
  \item The exclusion radius trades signal kept against noise removed and spurious signal; there is no radius that is best on every measure, and the choice depends on the transients of interest.
  \item On realistic fault slip, local filters keep the near-fault signal and remove the far field, which Gaussian transients and whole-event averages do not show.
  \item On the test years, no learned filter reduced the reconstruction error below that of the far stack, and each kept less signal; training with planted transients brought the attention network's signal kept up to, not beyond, that of the far stack.
  \item Scatter reduction, velocity repeatability and reconstruction error all reward filters that remove signal or the station's own noise; signal kept and spurious signal need to be reported alongside them.
\end{enumerate}

